\documentclass[11pt]{article}
\usepackage[margin=1in]{geometry}
\usepackage{amsmath,amssymb}
\usepackage{enumitem}
\usepackage{booktabs}
\usepackage{xcolor}
\usepackage{framed}
\usepackage{titlesec}

\titleformat{\section}{\large\bfseries}{}{0em}{}
\titleformat{\subsection}{\normalsize\bfseries}{}{0em}{}

\setlength{\parindent}{0pt}
\setlength{\parskip}{6pt}

\newcounter{qnum}
\newcommand{\question}[1]{\stepcounter{qnum}\subsection*{Question \theqnum \quad (#1)}}

\begin{document}

\begin{center}
{\LARGE\bfseries STAT GU4204 --- Practice Midterm Questions}\\[6pt]
{\large Covers all material through Lecture Slide Set 6}\\[4pt]
{\normalsize Question types: Fill-in-the-blank, True/False, Multiple Choice, Listing, Computation}
\end{center}

\bigskip
\hrule
\bigskip

% ====================================================================
% FILL IN THE BLANK
% ====================================================================
\section{Part A: Fill in the Blank}

\question{Fill in the Blank}
A function $T = \varphi(X_1, \ldots, X_n)$ that depends only on the observed data and contains no unknown parameters is called a \underline{\hspace{2.5in}}.

\question{Fill in the Blank}
An estimator $\hat\theta_n$ is said to be \underline{\hspace{2in}} if $\hat\theta_n \xrightarrow{P} \theta$ as $n \to \infty$.

\question{Fill in the Blank}
The decomposition $\text{MSE}(\hat\theta) = \text{Var}(\hat\theta) + [\text{bias}(\hat\theta)]^2$ is called the \underline{\hspace{2.5in}} decomposition.

\question{Fill in the Blank}
A statistic $T$ is \underline{\hspace{2in}} for $\theta$ if the conditional distribution of $\mathbf{X}$ given $T = t$ does not depend on $\theta$.

\question{Fill in the Blank}
The function $\dot\ell(x, \theta) = \frac{\partial}{\partial\theta}\log f(x, \theta)$ is called the \underline{\hspace{2in}}, and its expected value equals \underline{\hspace{1in}}.

\question{Fill in the Blank}
A random variable $\Psi(\mathbf{X}_n, g(\theta))$ whose distribution does not depend on $\theta$ is called a \underline{\hspace{2in}}.

\question{Fill in the Blank}
If a prior distribution and its resulting posterior distribution belong to the same parametric family, the prior is called a \underline{\hspace{2.5in}} prior.

\question{Fill in the Blank}
Under squared error loss, the Bayes estimator of $\theta$ is the \underline{\hspace{2.5in}} of the posterior distribution.

\question{Fill in the Blank}
If $X_1, \ldots, X_n$ are iid $N(\mu, \sigma^2)$, then $\frac{\sum_{i=1}^n (X_i - \bar X)^2}{\sigma^2}$ follows a \underline{\hspace{2in}} distribution.

\question{Fill in the Blank}
The property that if $\hat\theta$ is the MLE of $\theta$, then $g(\hat\theta)$ is the MLE of $g(\theta)$, is called \underline{\hspace{2in}} of the MLE.

\bigskip
\hrule
\bigskip

% ====================================================================
% TRUE / FALSE
% ====================================================================
\section{Part B: True or False}
\textit{Circle T or F. No justification required.}

\question{True/False}
\textbf{T \quad F} \quad If $Z_n \xrightarrow{d} Z$, then $Z_n \xrightarrow{P} Z$.

\question{True/False}
\textbf{T \quad F} \quad The MLE is always an unbiased estimator of the parameter.

\question{True/False}
\textbf{T \quad F} \quad If $\hat\theta_n$ is a consistent estimator of $\theta$, then $\hat\theta_n$ is unbiased for all $n$.

\question{True/False}
\textbf{T \quad F} \quad The Cram\'er-Rao lower bound applies to $X_1, \ldots, X_n$ iid Uniform$[0, \theta]$.

\question{True/False}
\textbf{T \quad F} \quad For $X_1, \ldots, X_n$ iid $N(\mu, \sigma^2)$, $\bar X$ and $s^2$ are independent random variables.

\question{True/False}
\textbf{T \quad F} \quad If $(2.3, 4.7)$ is a 95\% confidence interval for $\mu$, then $\mu$ lies in $(2.3, 4.7)$ with probability 0.95.

\question{True/False}
\textbf{T \quad F} \quad The CLT guarantees that $\bar X_n$ follows an exact normal distribution for all $n$.

\question{True/False}
\textbf{T \quad F} \quad A biased estimator can have lower MSE than any unbiased estimator of the same parameter.

\question{True/False}
\textbf{T \quad F} \quad If $T(\mathbf{X})$ is sufficient for $\theta$ and $\delta(\mathbf{X})$ is any estimator of $g(\theta)$, then $E[\delta(\mathbf{X}) | T]$ has MSE no larger than that of $\delta(\mathbf{X})$.

\question{True/False}
\textbf{T \quad F} \quad The $t_n$ distribution converges to $N(0,1)$ as $n \to \infty$.

\question{True/False}
\textbf{T \quad F} \quad For $X_1, \ldots, X_n$ iid from any distribution with finite variance, $\bar X_n$ is always the MVUE of $\mu$.

\question{True/False}
\textbf{T \quad F} \quad A sufficient statistic must be a single real number (cannot be vector-valued).

\question{True/False}
\textbf{T \quad F} \quad As the sample size $n \to \infty$, the posterior distribution concentrates around the true $\theta$ regardless of the prior.

\question{True/False}
\textbf{T \quad F} \quad For iid $\text{Uniform}[0,\theta]$ data, $2\bar X$ has lower MSE than $X_{(n)}$.

\bigskip
\hrule
\bigskip

% ====================================================================
% MULTIPLE CHOICE
% ====================================================================
\section{Part C: Multiple Choice}
\textit{Circle the best answer.}

\question{Multiple Choice}
Suppose $X_1, \ldots, X_n$ are iid $\text{Poisson}(\lambda)$. The MLE of $\lambda$ is:
\begin{enumerate}[label=(\alph*)]
\item $\frac{1}{n}\sum_{i=1}^n X_i^2$
\item $\bar X$
\item $\frac{1}{\bar X}$
\item $\sqrt{\bar X}$
\end{enumerate}

\question{Multiple Choice}
Which of the following is a sufficient statistic for $\theta$ when $X_1, \ldots, X_n$ are iid $\text{Exp}(\theta)$?
\begin{enumerate}[label=(\alph*)]
\item $\max(X_1, \ldots, X_n)$
\item $\min(X_1, \ldots, X_n)$
\item $\sum_{i=1}^n X_i$
\item $\bar X^2$
\end{enumerate}

\question{Multiple Choice}
Suppose $X_1, \ldots, X_{25}$ are iid $N(\mu, 9)$ with $\sigma^2 = 9$ known. A 95\% CI for $\mu$ is:
\begin{enumerate}[label=(\alph*)]
\item $\bar X \pm 1.96 \cdot \frac{3}{25}$
\item $\bar X \pm 1.96 \cdot \frac{3}{\sqrt{25}}$
\item $\bar X \pm 1.96 \cdot \frac{9}{\sqrt{25}}$
\item $\bar X \pm t_{24, 0.025} \cdot \frac{3}{\sqrt{25}}$
\end{enumerate}

\question{Multiple Choice}
The Fisher information $I(\theta)$ for a single observation from $\text{Bernoulli}(\theta)$ is:
\begin{enumerate}[label=(\alph*)]
\item $\theta(1-\theta)$
\item $\frac{1}{\theta}$
\item $\frac{1}{\theta(1-\theta)}$
\item $\frac{1}{\theta^2}$
\end{enumerate}

\question{Multiple Choice}
If $Z \sim N(0,1)$ and $V \sim \chi^2_n$ are independent, then $T = Z / \sqrt{V/n}$ follows:
\begin{enumerate}[label=(\alph*)]
\item $N(0, 1)$
\item $\chi^2_{n+1}$
\item $t_n$
\item $F_{1,n}$
\end{enumerate}

\question{Multiple Choice}
The Cram\'er-Rao inequality does NOT apply when:
\begin{enumerate}[label=(\alph*)]
\item $X_i$ iid $N(\mu, \sigma^2)$
\item $X_i$ iid $\text{Poisson}(\lambda)$
\item $X_i$ iid $\text{Uniform}[0, \theta]$
\item $X_i$ iid $\text{Exponential}(\theta)$
\end{enumerate}

\question{Multiple Choice}
For Bayesian inference with $X_1, \ldots, X_n$ iid $\text{Bernoulli}(\theta)$ and prior $\theta \sim \text{Beta}(\alpha, \beta)$, the posterior distribution of $\theta$ is:
\begin{enumerate}[label=(\alph*)]
\item $\text{Beta}(\alpha + n, \beta)$
\item $\text{Beta}(\alpha + \sum x_i, \beta + n - \sum x_i)$
\item $\text{Gamma}(\alpha + \sum x_i, \beta + n)$
\item $N(\bar x, 1/n)$
\end{enumerate}

\bigskip
\hrule
\bigskip

% ====================================================================
% LISTING / SHORT ANSWER
% ====================================================================
\section{Part D: Listing / Short Answer}

\question{Listing}
List the three regularity conditions (A.1, A.2, A.3) required for the Cram\'er-Rao inequality.

\vspace{2.5in}

\question{Listing}
List three properties of the MLE. For each, state whether it always holds or requires conditions.

\vspace{2in}

\question{Short Answer}
State the Weak Law of Large Numbers. What are the conditions required?

\vspace{1.5in}

\question{Short Answer}
State the Central Limit Theorem precisely. What are the conditions?

\vspace{1.5in}

\question{Short Answer}
State the Factorization Criterion for sufficient statistics.

\vspace{1.5in}

\question{Short Answer}
State Proposition 6.3 (the three results about $\bar X$ and $s^2$ for a normal random sample).

\vspace{1.5in}

\question{Listing}
Name two estimation methods covered in this course and briefly describe each (one sentence each).

\vspace{1.5in}

\question{Short Answer}
Explain the difference between an \emph{estimator} and an \emph{estimate}. Give an example.

\vspace{1.5in}

\question{Short Answer}
What does it mean for a confidence interval to have ``confidence level $1 - \alpha$''? Why is it incorrect to say ``$\theta$ is in the interval with probability $1 - \alpha$''?

\vspace{1.5in}

\bigskip
\hrule
\bigskip

% ====================================================================
% COMPUTATIONAL
% ====================================================================
\section{Part E: Computational Problems}

\question{MLE Computation}
Suppose $X_1, \ldots, X_n$ are iid $\text{Exponential}(\theta)$ with pdf $f(x|\theta) = \theta e^{-\theta x}$ for $x > 0$.

\begin{enumerate}[label=(\alph*)]
\item Write down the likelihood function $L(\theta)$ and the log-likelihood $\ell(\theta)$.
\item Find the MLE $\hat\theta$ of $\theta$.
\item Is $\hat\theta$ unbiased? Justify.
\item Using the invariance property, find the MLE of $E(X_1) = 1/\theta$.
\end{enumerate}

\vspace{3in}

\question{Method of Moments}
Suppose $X_1, \ldots, X_n$ are iid $\text{Gamma}(\alpha, \beta)$ with $E(X) = \alpha/\beta$ and $\text{Var}(X) = \alpha/\beta^2$.

\begin{enumerate}[label=(\alph*)]
\item Find the method of moments estimators $\hat\alpha$ and $\hat\beta$.
\item Are these estimators consistent? Explain why or why not.
\end{enumerate}

\vspace{2.5in}

\question{Sufficient Statistics}
Let $X_1, \ldots, X_n$ be iid $N(\mu, \sigma^2)$ where both $\mu$ and $\sigma^2$ are unknown.

\begin{enumerate}[label=(\alph*)]
\item Write the joint pdf $f(\mathbf{x} | \mu, \sigma^2)$.
\item Use the Factorization Criterion to find a sufficient statistic for $(\mu, \sigma^2)$.
\item Why is a single statistic (like $\sum X_i$) not sufficient in this case?
\end{enumerate}

\vspace{3in}

\question{MSE Comparison}
Suppose $X_1, \ldots, X_n$ are iid $\text{Uniform}[0, \theta]$.

Consider two estimators of $\theta$:
\begin{itemize}
\item $T_1 = 2\bar X$
\item $T_2 = \left(1 + \frac{1}{n}\right) X_{(n)}$
\end{itemize}

\begin{enumerate}[label=(\alph*)]
\item Show that both $T_1$ and $T_2$ are unbiased for $\theta$.

\textit{Hint:} $E(X_{(n)}) = \frac{n}{n+1}\theta$.

\item Compute $\text{MSE}(T_1)$.

\textit{Hint:} $\text{Var}(X_i) = \theta^2/12$.

\item Given that $\text{MSE}(T_2) = \frac{\theta^2}{n(n+2)}$, which estimator is better for $n = 5$? For $n = 50$?
\end{enumerate}

\vspace{3in}

\question{Delta Method}
Let $X_1, \ldots, X_n$ be iid $\text{Bernoulli}(\theta)$ where $0 < \theta < 1$.

\begin{enumerate}[label=(\alph*)]
\item State the asymptotic distribution of $\bar X_n$ (using the CLT).
\item Using the delta method, find the asymptotic distribution of $g(\bar X_n) = \log\!\left(\frac{\bar X_n}{1 - \bar X_n}\right)$ (the log-odds).

\textit{Hint:} $g'(x) = \frac{1}{x(1-x)}$.
\end{enumerate}

\vspace{3in}

\question{Confidence Interval Construction}
Suppose $X_1, \ldots, X_{16}$ are iid $N(\mu, \sigma^2)$ with $\sigma^2$ \textbf{unknown}. You observe $\bar x = 5.2$ and $s = 2.4$.

\begin{enumerate}[label=(\alph*)]
\item What pivot do you use to construct a CI for $\mu$? State its distribution.
\item Construct a 95\% confidence interval for $\mu$.

\textit{You may use} $t_{15, 0.025} = 2.131$.

\item How would your answer change if $\sigma = 2.4$ were \emph{known} instead? Which interval is wider and why?
\end{enumerate}

\vspace{3in}

\question{Fisher Information and CRLB}
Let $X_1, \ldots, X_n$ be iid $\text{Poisson}(\theta)$ with pmf $f(x|\theta) = e^{-\theta}\theta^x / x!$.

\begin{enumerate}[label=(\alph*)]
\item Compute the score function $\dot\ell(x, \theta)$.
\item Compute the Fisher information $I(\theta)$.
\item State the Cram\'er-Rao lower bound for any unbiased estimator of $\theta$.
\item Is $\bar X$ the MVUE of $\theta$? Justify.
\end{enumerate}

\vspace{3in}

\question{Bayesian Inference}
Suppose $X_1, \ldots, X_n$ are iid $\text{Poisson}(\lambda)$ and we use the prior $\lambda \sim \text{Gamma}(\alpha, \beta)$.

\begin{enumerate}[label=(\alph*)]
\item Find the posterior distribution of $\lambda | \mathbf{x}$.
\item Find the Bayes estimator of $\lambda$ under squared error loss.
\item Show that as $n \to \infty$, the Bayes estimator converges to the MLE.
\end{enumerate}

\vspace{3in}

\question{Sampling Distributions}
Suppose $X_1, \ldots, X_9$ are iid $N(50, 16)$.

\begin{enumerate}[label=(\alph*)]
\item What is the distribution of $\bar X$?
\item What is the distribution of $\frac{\sum_{i=1}^9 (X_i - \bar X)^2}{16}$?
\item What is the distribution of $\frac{\sqrt{9}(\bar X - 50)}{s}$?
\item Compute $P(\bar X > 52)$.
\end{enumerate}

\vspace{3in}

\question{Rao-Blackwell}
Let $X_1, \ldots, X_n$ be iid $\text{Bernoulli}(\theta)$.

\begin{enumerate}[label=(\alph*)]
\item Show that $T = \sum_{i=1}^n X_i$ is sufficient for $\theta$.
\item Consider the crude estimator $\delta(X_1, \ldots, X_n) = X_1$ (just the first observation). Show it is unbiased for $\theta$.
\item Compute $E[X_1 | T = t]$ and verify it improves upon $X_1$.

\textit{Hint:} Given $T = t$, each $X_i$ is equally likely to be 1 or 0, so $E[X_1 | T = t] = t/n$.
\end{enumerate}

\vspace{3in}

\question{Chi-Squared and CI for Variance}
Suppose $X_1, \ldots, X_{21}$ are iid $N(\mu, \sigma^2)$ and we observe $s^2 = 12.5$.

\begin{enumerate}[label=(\alph*)]
\item What is the distribution of $(n-1)s^2/\sigma^2$?
\item Construct a 90\% confidence interval for $\sigma^2$.

\textit{You may use:} $\chi^2_{20, 0.05} = 31.41$ and $\chi^2_{20, 0.95} = 10.85$.
\end{enumerate}

\vspace{2.5in}

\question{Chebyshev / Sample Size}
A machine fills bottles with a mean of $\mu$ ounces and a standard deviation of $\sigma = 0.5$ ounces. You want to estimate $\mu$ by $\bar X_n$.

\begin{enumerate}[label=(\alph*)]
\item Using Chebyshev's inequality, how large must $n$ be so that $P(|\bar X_n - \mu| \ge 0.1) \le 0.01$?
\item Using the CLT approximation instead, how large must $n$ be so that a 99\% CI for $\mu$ has total width at most 0.2?

\textit{You may use} $z_{0.005} = 2.576$.
\end{enumerate}

\vspace{2.5in}

\question{Consistency}
Let $X_1, \ldots, X_n$ be iid with mean $\mu$ and variance $\sigma^2 < \infty$.

\begin{enumerate}[label=(\alph*)]
\item Prove that $\bar X_n$ is a consistent estimator of $\mu$ using Chebyshev's inequality.
\item Using the continuous mapping theorem, show that $\bar X_n^2$ is a consistent estimator of $\mu^2$.
\end{enumerate}

\vspace{2.5in}

\question{MLE for Uniform}
Suppose $X_1, \ldots, X_n$ are iid $\text{Uniform}[0, \theta]$.

\begin{enumerate}[label=(\alph*)]
\item Write the likelihood $L(\theta)$.
\item Explain why you cannot set $L'(\theta) = 0$ to find the MLE.
\item Find the MLE of $\theta$ by reasoning about the shape of $L(\theta)$.
\item Is the MLE unbiased? If not, find its bias.

\textit{Hint:} $E(X_{(n)}) = \frac{n}{n+1}\theta$.
\end{enumerate}

\vspace{3in}

\question{Normal MLE Bias}
Let $X_1, \ldots, X_n$ be iid $N(\mu, \sigma^2)$.

\begin{enumerate}[label=(\alph*)]
\item The MLE of $\sigma^2$ is $\hat\sigma^2 = \frac{1}{n}\sum_{i=1}^n(X_i - \bar X)^2$. Show that $E(\hat\sigma^2) = \frac{n-1}{n}\sigma^2$.

\textit{Hint:} Use the fact that $\sum(X_i - \bar X)^2/\sigma^2 \sim \chi^2_{n-1}$.

\item Is $\hat\sigma^2$ unbiased? What modification makes it unbiased?
\item For $n = 10$, compute the bias of $\hat\sigma^2$.
\end{enumerate}

\vspace{3in}

\question{CRLB with $g(\theta)$}
Let $X_1, \ldots, X_n$ be iid $\text{Bernoulli}(\theta)$. Suppose we want to estimate $g(\theta) = \theta^2$.

\begin{enumerate}[label=(\alph*)]
\item Compute $g'(\theta)$.
\item Using the Fisher information $I(\theta) = \frac{1}{\theta(1-\theta)}$, state the CRLB for any unbiased estimator of $\theta^2$.
\item Simplify the bound. Does $(\bar X)^2$ achieve this bound?
\end{enumerate}

\vspace{2.5in}

\bigskip
\hrule
\bigskip

% ====================================================================
% ANSWERS
% ====================================================================
\newpage
\section*{Answer Key}

\subsection*{Part A: Fill in the Blank}
\begin{enumerate}
\item statistic
\item consistent
\item bias-variance
\item sufficient
\item score function; zero (0)
\item pivot
\item conjugate
\item mean (posterior mean)
\item $\chi^2_{n-1}$ (chi-squared with $n-1$ degrees of freedom)
\item invariance
\end{enumerate}

\subsection*{Part B: True or False}
\begin{enumerate}[start=11]
\item \textbf{F} --- conv in dist does NOT imply conv in prob (only if limit is constant)
\item \textbf{F} --- MLE of $\sigma^2$ in Normal is biased; $X_{(n)}$ for Uniform is biased
\item \textbf{F} --- consistency is a large-$n$ property; estimator can be biased for each finite $n$
\item \textbf{F} --- Uniform$[0,\theta]$ violates regularity condition A.1 (support depends on $\theta$)
\item \textbf{T} --- Proposition 6.3; this independence characterizes the normal distribution
\item \textbf{F} --- $\theta$ is fixed; the probability statement is about the random interval before data is observed
\item \textbf{F} --- CLT gives an approximate distribution for large $n$; exact normality requires $X_i$ to be normal
\item \textbf{T} --- e.g., $X_{(n)}$ for Uniform$[0,\theta]$ has lower MSE than any unbiased estimator like $2\bar X$
\item \textbf{T} --- Rao-Blackwell theorem
\item \textbf{T} --- by WLLN, $V/n \xrightarrow{P} 1$, so denominator $\to 1$ and $T \xrightarrow{d} N(0,1)$
\item \textbf{F} --- $\bar X_n$ is unbiased for $\mu$ always, but MVUE depends on the model (need CRLB or other argument)
\item \textbf{F} --- e.g., $(\sum X_i, \sum X_i^2)$ is sufficient for $N(\mu,\sigma^2)$ with both unknown
\item \textbf{T} --- posterior consistency; data overwhelms the prior
\item \textbf{F} --- $\text{MSE}(2\bar X) = \theta^2/(3n)$ while $\text{MSE}(X_{(n)})$ is of order $\theta^2/n^2$, much smaller
\end{enumerate}

\subsection*{Part C: Multiple Choice}
\begin{enumerate}[start=25]
\item \textbf{(b)} $\bar X$
\item \textbf{(c)} $\sum X_i$
\item \textbf{(b)} $\bar X \pm 1.96 \cdot 3/\sqrt{25} = \bar X \pm 1.176$
\item \textbf{(c)} $I(\theta) = 1/[\theta(1-\theta)]$
\item \textbf{(c)} $t_n$
\item \textbf{(c)} Uniform$[0,\theta]$ --- support depends on $\theta$, violating A.1
\item \textbf{(b)} Beta$(\alpha + \sum x_i, \;\beta + n - \sum x_i)$
\end{enumerate}

\subsection*{Part D: Listing / Short Answer}
\begin{enumerate}[start=32]
\item \textbf{A.1}: Support $\{x: f(x,\theta)>0\}$ does not depend on $\theta$. \textbf{A.2}: Can interchange $\partial/\partial\theta$ and $\int$. \textbf{A.3}: $\partial/\partial\theta \log f(x,\theta)$ exists and is finite for all $x, \theta$.

\item \textbf{Invariance} (always holds): $g(\hat\theta)$ is MLE of $g(\theta)$. \textbf{Consistency} (under regularity conditions): $\hat\theta_n \xrightarrow{P} \theta$. \textbf{May be biased / may not exist / may not be unique} (examples: Normal $\hat\sigma^2$ is biased; Uniform$[\theta,\theta+1]$ MLE is not unique).

\item WLLN: If $X_1,\ldots,X_n$ are iid with $E|X_1| < \infty$, then $\bar X_n \xrightarrow{P} E(X_1) = \mu$. Conditions: iid, finite mean.

\item CLT: If $X_i$ iid with mean $\mu$ and $0 < \sigma^2 < \infty$, then $\sqrt{n}(\bar X_n - \mu)/\sigma \xrightarrow{d} N(0,1)$. Conditions: iid, finite variance.

\item Factorization: $T = r(\mathbf{X})$ is sufficient for $\theta$ iff $f(\mathbf{x},\theta) = u(\mathbf{x}) \cdot v(r(\mathbf{x}), \theta)$ where $u$ depends only on data and $v$ depends on $\theta$ only through $r(\mathbf{x})$.

\item Prop 6.3: For $X_i$ iid $N(\mu,\sigma^2)$: (1) $\bar X \sim N(\mu, \sigma^2/n)$; (2) $s^2 \sim \frac{\sigma^2}{n-1}\chi^2_{n-1}$; (3) $\bar X$ and $s^2$ are independent.

\item MOM: Set sample moments equal to population moments, solve. MLE: Find $\theta$ maximizing the likelihood $L(\theta) = \prod f(X_i, \theta)$.

\item An estimator is a random variable (function of data); an estimate is the realized number. E.g., $\hat\mu = \bar X_n$ is the estimator; $\hat\mu = 3.7$ is an estimate.

\item A level $1{-}\alpha$ CI means that before observing data, the random interval contains $g(\theta)$ with probability $\ge 1{-}\alpha$. After data is observed, $\theta$ is fixed and the realized interval either contains it or doesn't --- we cannot assign a probability.
\end{enumerate}

\subsection*{Part E: Computational Problems}

\begin{enumerate}[start=41]
\item \textbf{(a)} $L(\theta) = \theta^n e^{-\theta \sum x_i}$, $\ell(\theta) = n\log\theta - \theta\sum x_i$. \textbf{(b)} $\ell'(\theta) = n/\theta - \sum x_i = 0 \Rightarrow \hat\theta = n/\sum x_i = 1/\bar X$. \textbf{(c)} $E(1/\bar X) \ne \theta$ in general (Jensen's inequality), so biased. \textbf{(d)} By invariance, MLE of $1/\theta$ is $1/\hat\theta = \bar X$.

\item \textbf{(a)} $E(X) = \alpha/\beta = \bar X$ and $E(X^2) - [E(X)]^2 = \alpha/\beta^2$, so $\alpha/\beta^2 = \hat\mu_2 - \bar X^2$. Solving: $\hat\beta = \bar X/(\hat\mu_2 - \bar X^2)$ and $\hat\alpha = \bar X^2/(\hat\mu_2 - \bar X^2)$. \textbf{(b)} Yes, by LLN $\bar X \xrightarrow{P} \alpha/\beta$ and $\hat\mu_2 \xrightarrow{P} E(X^2)$, then CMT.

\item \textbf{(a)} $f(\mathbf{x}) = (2\pi\sigma^2)^{-n/2}\exp\left(-\frac{1}{2\sigma^2}\sum(x_i - \mu)^2\right)$. \textbf{(b)} Expand: $\sum(x_i-\mu)^2 = \sum x_i^2 - 2\mu\sum x_i + n\mu^2$. So $f = \underbrace{1}_{u(\mathbf{x})} \cdot (2\pi\sigma^2)^{-n/2}\exp\left(-\frac{\sum x_i^2 - 2\mu\sum x_i + n\mu^2}{2\sigma^2}\right)$. Depends on data through $(\sum x_i, \sum x_i^2)$. \textbf{(c)} Two unknown parameters require (at least) two-dimensional sufficient statistic.

\item \textbf{(a)} $E(T_1) = 2\mu = 2\cdot\theta/2 = \theta$. $E(T_2) = (1+1/n)\cdot\frac{n}{n+1}\theta = \theta$. Both unbiased. \textbf{(b)} $\text{Var}(T_1) = 4\text{Var}(\bar X) = 4\cdot\theta^2/(12n) = \theta^2/(3n)$. Since unbiased, $\text{MSE} = \theta^2/(3n)$. \textbf{(c)} $n=5$: $T_1$ MSE $= \theta^2/15 \approx 0.067\theta^2$; $T_2$ MSE $= \theta^2/35 \approx 0.029\theta^2$. $T_2$ better. $n=50$: $T_1$ MSE $= \theta^2/150$; $T_2$ MSE $= \theta^2/2600$. $T_2$ much better.

\item \textbf{(a)} $\sqrt{n}(\bar X_n - \theta)/\sqrt{\theta(1-\theta)} \xrightarrow{d} N(0,1)$, so $\bar X_n \sim_{\text{approx}} N(\theta, \theta(1-\theta)/n)$. \textbf{(b)} $g'(\theta) = 1/[\theta(1-\theta)]$. By delta method: $\sqrt{n}(g(\bar X_n) - g(\theta)) \xrightarrow{d} N\!\left(0,\; \frac{\theta(1-\theta)}{1} \cdot \frac{1}{\theta^2(1-\theta)^2}\right) = N\!\left(0, \frac{1}{\theta(1-\theta)}\right)$.

\item \textbf{(a)} Pivot: $\sqrt{n}(\bar X - \mu)/s \sim t_{n-1} = t_{15}$. \textbf{(b)} $5.2 \pm 2.131 \cdot 2.4/\sqrt{16} = 5.2 \pm 2.131 \cdot 0.6 = 5.2 \pm 1.279 = (3.921, 6.479)$. \textbf{(c)} If $\sigma$ known: $5.2 \pm 1.96 \cdot 0.6 = 5.2 \pm 1.176 = (4.024, 6.376)$. The $t$-interval is wider because $t_{15,0.025} = 2.131 > z_{0.025} = 1.96$ (heavier tails compensate for estimating $\sigma$).

\item \textbf{(a)} $\ell(x,\theta) = -\theta + x\log\theta - \log(x!)$. $\dot\ell = -1 + x/\theta = (x-\theta)/\theta$. \textbf{(b)} $I(\theta) = E[(X-\theta)^2/\theta^2] = \text{Var}(X)/\theta^2 = \theta/\theta^2 = 1/\theta$. \textbf{(c)} CRLB: $\text{Var}(T) \ge 1/(n\cdot 1/\theta) = \theta/n$. \textbf{(d)} $\text{Var}(\bar X) = \theta/n = $ CRLB. Since $\bar X$ is unbiased and achieves the bound, it is the MVUE.

\item \textbf{(a)} Posterior $\propto \prod \frac{e^{-\lambda}\lambda^{x_i}}{x_i!} \cdot \frac{\beta^\alpha}{\Gamma(\alpha)}\lambda^{\alpha-1}e^{-\beta\lambda} \propto \lambda^{\alpha + \sum x_i - 1}e^{-(\beta+n)\lambda}$. This is Gamma$(\alpha + \sum x_i, \;\beta + n)$. \textbf{(b)} Posterior mean: $\frac{\alpha + \sum x_i}{\beta + n}$. \textbf{(c)} As $n \to \infty$: $\frac{\alpha + \sum x_i}{\beta + n} = \frac{\alpha/n + \bar x}{(\beta/n) + 1} \to \bar x = $ MLE.

\item \textbf{(a)} $\bar X \sim N(50, 16/9)$, i.e., $N(50, 16/9)$. \textbf{(b)} $\chi^2_8$. \textbf{(c)} $t_8$. \textbf{(d)} $P(\bar X > 52) = P\!\left(Z > \frac{52-50}{4/3}\right) = P(Z > 1.5) = 1 - \Phi(1.5) \approx 0.0668$.

\item \textbf{(a)} Joint pmf: $\prod \theta^{x_i}(1-\theta)^{1-x_i} = \theta^{\sum x_i}(1-\theta)^{n - \sum x_i} = \underbrace{1}_{u(\mathbf{x})}\cdot v(\sum x_i, \theta)$. So $T = \sum X_i$ is sufficient. \textbf{(b)} $E(X_1) = \theta$, so unbiased. \textbf{(c)} $E[X_1|T=t] = t/n = \bar X$. MSE$(\bar X) = \theta(1-\theta)/n \le \theta(1-\theta) = $ MSE$(X_1)$. Strict improvement for $n \ge 2$.

\item \textbf{(a)} $(n-1)s^2/\sigma^2 = 20s^2/\sigma^2 \sim \chi^2_{20}$. \textbf{(b)} $P\!\left(\chi^2_{20,0.95} \le \frac{20s^2}{\sigma^2} \le \chi^2_{20,0.05}\right) = 0.90$. Inverting: $\sigma^2 \in \left(\frac{20 \cdot 12.5}{31.41},\; \frac{20 \cdot 12.5}{10.85}\right) = (7.96, 23.04)$.

\item \textbf{(a)} Need $\sigma^2/(n\varepsilon^2) \le 0.01$: $0.25/(n \cdot 0.01) \le 0.01 \Rightarrow n \ge 2500$. \textbf{(b)} Width $= 2 \cdot z_{0.005}\cdot\sigma/\sqrt{n} \le 0.2$: $2(2.576)(0.5)/\sqrt{n} \le 0.2 \Rightarrow \sqrt{n} \ge 12.88 \Rightarrow n \ge 166$.

\item \textbf{(a)} $P(|\bar X_n - \mu| > \varepsilon) \le \text{Var}(\bar X_n)/\varepsilon^2 = \sigma^2/(n\varepsilon^2) \to 0$. \textbf{(b)} $g(x) = x^2$ is continuous. By CMT, $\bar X_n \xrightarrow{P} \mu \Rightarrow \bar X_n^2 \xrightarrow{P} \mu^2$.

\item \textbf{(a)} $L(\theta) = \prod_{i=1}^n \frac{1}{\theta}I(0 \le x_i \le \theta) = \theta^{-n} \cdot I(\theta \ge x_{(n)})$. \textbf{(b)} $L$ is not differentiable at $\theta = x_{(n)}$ (discontinuity); for $\theta > x_{(n)}$, $L'(\theta) = -n\theta^{-(n+1)} < 0$ (always decreasing). \textbf{(c)} $L$ is maximized at the smallest allowable $\theta$, which is $\hat\theta = X_{(n)}$. \textbf{(d)} $E(\hat\theta) = \frac{n}{n+1}\theta \ne \theta$. Bias $= -\frac{\theta}{n+1}$ (underestimates $\theta$).

\item \textbf{(a)} $E(\hat\sigma^2) = \frac{1}{n}E\!\left[\sum(X_i-\bar X)^2\right] = \frac{\sigma^2}{n}E(\chi^2_{n-1}) = \frac{\sigma^2(n-1)}{n}$. \textbf{(b)} Biased; $E(\hat\sigma^2) = \frac{n-1}{n}\sigma^2 \ne \sigma^2$. Modification: $s^2 = \frac{n}{n-1}\hat\sigma^2 = \frac{1}{n-1}\sum(X_i-\bar X)^2$. \textbf{(c)} For $n=10$: bias $= \frac{9}{10}\sigma^2 - \sigma^2 = -\sigma^2/10$.

\item \textbf{(a)} $g'(\theta) = 2\theta$. \textbf{(b)} CRLB $= \frac{(2\theta)^2}{n \cdot 1/[\theta(1-\theta)]} = \frac{4\theta^2 \cdot \theta(1-\theta)}{n} = \frac{4\theta^3(1-\theta)}{n}$. \textbf{(c)} $(\bar X)^2$ is biased for $\theta^2$ ($E(\bar X^2) = \theta^2 + \theta(1-\theta)/n \ne \theta^2$), so CRLB does not directly apply to it (it's not unbiased).
\end{enumerate}

\bigskip
\hrule
\bigskip

% ====================================================================
% WORKFLOW DRILL PROBLEMS
% ====================================================================
\section{Part F: Workflow Drill Problems}
\textit{Each problem is tagged with the workflow type [W1--W10] from the cheat sheet workflow guide.}

\setcounter{qnum}{0}

\question{Variance of Linear Combinations [W1]}
Let $X, Y, Z$ be random variables with $\text{Var}(X)=5$, $\text{Var}(Y)=3$, $\text{Var}(Z)=2$, $\text{Cov}(X,Y)=1$, $\text{Cov}(X,Z)=-2$, $\text{Cov}(Y,Z)=0$.

\begin{enumerate}[label=(\alph*)]
\item Find $\text{Var}(2X - Y + 3Z)$.
\item Find $\text{Var}(X + Y - Z + 7)$.
\end{enumerate}

\question{Chebyshev Bound \& Sample Size [W2]}
Let $X_1, \ldots, X_n$ be iid with mean $\mu = 10$ and variance $\sigma^2 = 25$.

\begin{enumerate}[label=(\alph*)]
\item Use Chebyshev's inequality to find an upper bound on $P(|\bar{X}_n - 10| \geq 3)$.
\item How large must $n$ be so that $P(|\bar{X}_n - 10| \leq 1) \geq 0.95$?
\end{enumerate}

\question{Markov Inequality [W2]}
Let $X$ be a non-negative random variable with $E(X) = 4$.

\begin{enumerate}[label=(\alph*)]
\item Find an upper bound on $P(X \geq 20)$.
\item If additionally $E(|X - E(X)|^4) = 1000$, find an upper bound on $P(|X - E(X)| \geq 5)$ using a 4th-moment bound.
\end{enumerate}

\question{Standard Delta Method [W3]}
Let $X_1, \ldots, X_n$ be iid with mean $\theta$ and variance $\sigma^2$. Find the asymptotic distribution of $(\bar{X}_n)^4$.

\question{Delta Method --- Transformed Variance [W3]}
Let $Y_n = \frac{1}{n}\sum_{i=1}^n X_i^2$, where $E(X^2) = \sigma^2$ and $\text{Var}(X^2) = 2\sigma^4$. Find the asymptotic distribution of $\sqrt{Y_n}$.

\question{Order Statistics [W4]}
Let $X_1, \ldots, X_n$ be iid Unif$[0, \theta]$. Let $Y_n = \max(X_i)$.

\begin{enumerate}[label=(\alph*)]
\item Find the CDF of $Y_n$.
\item Find the CDF of $Z_n = n(Y_n - \theta)$ and determine its limiting distribution as $n \to \infty$.
\end{enumerate}

\question{MLE --- Calculus Path [W5]}
Let $X_1, \ldots, X_n$ be iid with PDF $f(x|\theta) = \theta x^{\theta-1}$ for $0 < x < 1$, $\theta > 0$. Find the MLE of $\theta$.

\question{MLE --- Restricted Domain [W5]}
Let $X_1, \ldots, X_{70}$ be iid Bernoulli$(p)$. You observe $\sum_{i=1}^{70} x_i = 58$. Find the MLE of $p$ if the parameter space is restricted to $p \in [1/2, 2/3]$.

\question{MLE Invariance [W6]}
Let $X_1, \ldots, X_n$ be iid Exponential$(\beta)$. The MLE of $\beta$ is $\hat\beta = 1/\bar{X}$.

\begin{enumerate}[label=(\alph*)]
\item Find the MLE of the median of the distribution.
\item Find the MLE of $P(X > 2)$.
\end{enumerate}

\question{MOM Estimator [W7]}
Let $X_1, \ldots, X_n$ be iid Bernoulli$(\theta)$. Find the method of moments estimator of $\theta$ and verify it equals the MLE.

\question{Consistency Proof [W8]}
Let $X_1, \ldots, X_n$ be iid Exponential$(\beta)$ with mean $1/\beta$. Show that $\hat\beta = 1/\bar{X}_n$ is consistent for $\beta$.

\question{Gamma Integral [W9]}
Evaluate $\displaystyle\int_0^\infty x^5 e^{-3x}\,dx$.

\question{Expected Power of Gamma [W9]}
Let $X \sim \text{Gamma}(2, 3)$. Compute $E[X^4]$.

\question{Sum of Exponentials [W10]}
Let $X_1, \ldots, X_8$ be iid Exponential$(\beta)$. Find the distribution of $S = X_1 + \cdots + X_8$.

\question{Building a $t$-statistic [W10]}
Let $X_1, \ldots, X_5$ be iid $N(0,1)$. Define $W = X_1 / \sqrt{(X_2^2 + X_3^2 + X_4^2)/3}$.

\begin{enumerate}[label=(\alph*)]
\item What is the distribution of $W$?
\item What is the distribution of $W^2$?
\end{enumerate}

\bigskip
\hrule
\bigskip

% ====================================================================
% WORKFLOW DRILL ANSWERS
% ====================================================================
\section{Workflow Drill Answers}

\begin{enumerate}
\item \textbf{(a)} $4(5) + 1(3) + 9(2) + 2(2)(-1)(1) + 2(2)(3)(-2) + 2(-1)(3)(0) = 20 + 3 + 18 - 4 - 24 + 0 = \boxed{13}$. \textbf{(b)} $1(5) + 1(3) + 1(2) + 2(1)(1)(1) + 2(1)(-1)(-2) + 2(1)(-1)(0) = 5 + 3 + 2 + 2 + 4 + 0 = \boxed{16}$

\item \textbf{(a)} $P(|\bar X - 10| \geq 3) \leq \text{Var}(\bar X)/9 = 25/(9n) = \boxed{25/(9n)}$. \textbf{(b)} Need $25/(n \cdot 1) \leq 0.05$, so $\boxed{n \geq 500}$.

\item \textbf{(a)} Markov: $P(X \geq 20) \leq E(X)/20 = 4/20 = \boxed{1/5}$. \textbf{(b)} $P(|X - \mu| \geq 5) \leq \beta_4/5^4 = 1000/625 = 1.6$. Bound exceeds 1, so trivially $\leq 1$ (uninformative).

\item $g(x) = x^4$, $g'(\theta) = 4\theta^3$. By delta method: $\boxed{N(\theta^4,\; 16\sigma^2\theta^6/n)}$.

\item $g(x) = \sqrt{x}$, $g'(\sigma^2) = 1/(2\sigma)$. Asymptotic variance: $(2\sigma^4/n) \cdot 1/(4\sigma^2) = \sigma^2/(2n)$. Answer: $\boxed{N(\sigma,\; \sigma^2/(2n))}$.

\item \textbf{(a)} $P(Y_n \leq y) = (y/\theta)^n$ for $0 < y < \theta$. \textbf{(b)} $P(Z_n \leq z) = (1 + z/(n\theta))^n \to \boxed{e^{z/\theta}}$ for $z \leq 0$.

\item $\ell(\theta) = n\log\theta + (\theta - 1)\sum\log x_i$. Setting $\ell'(\theta) = n/\theta + \sum\log x_i = 0$: $\boxed{\hat\theta = -n/\sum\log x_i}$. Check: $\ell'' = -n/\theta^2 < 0$. Note $\hat\theta > 0$ since $\log x_i < 0$.

\item Unrestricted MLE: $\hat p = 58/70 \approx 0.829 > 2/3$. Falls outside domain. $L(p) = p^{58}(1-p)^{12}$ is increasing on $[1/2, 2/3]$ (since maximum is to the right). So $\boxed{\hat p = 2/3}$.

\item \textbf{(a)} Median of Exp$(\beta)$ is $(\log 2)/\beta$. By MLE invariance: $\hat m = (\log 2)/\hat\beta = \boxed{(\log 2)\cdot\bar X}$. \textbf{(b)} $P(X > 2) = e^{-2\beta}$. By invariance: $\boxed{e^{-2/\bar X}}$.

\item MOM: $E(X) = \theta$, so $\hat\theta_{\text{MOM}} = \bar X$. MLE: $\ell = \sum x_i\log\theta + (n - \sum x_i)\log(1-\theta)$. Setting $\ell' = 0$: $\hat\theta_{\text{MLE}} = \bar X$. $\boxed{\text{MOM} = \text{MLE} = \bar X}$.

\item By WLLN: $\bar X_n \xrightarrow{P} E(X) = 1/\beta$. Let $g(x) = 1/x$; $g$ is continuous at $1/\beta > 0$. By CMT: $\boxed{1/\bar X_n \xrightarrow{P} \beta}$.

\item Match to $\int_0^\infty x^{\alpha-1}e^{-\beta x}dx$ with $\alpha = 6$, $\beta = 3$. Answer: $\Gamma(6)/3^6 = 120/729 = \boxed{40/243}$.

\item $E[X^4] = \Gamma(2+4)/\!\left(3^4 \cdot \Gamma(2)\right) = \Gamma(6)/\!\left(81 \cdot 1\right) = 120/81 = \boxed{40/27}$.

\item MGF of $S$: $\prod_{i=1}^8 \frac{\beta}{\beta - t} = \left(\frac{\beta}{\beta - t}\right)^8$. This is the MGF of $\boxed{\text{Gamma}(8, \beta)}$.

\item \textbf{(a)} $X_1 \sim N(0,1)$ and $V = X_2^2 + X_3^2 + X_4^2 \sim \chi^2_3$, independent of $X_1$. So $W = X_1/\sqrt{V/3} \sim \boxed{t_3}$. \textbf{(b)} $W^2$ has an $F(1,3)$ distribution (square of a $t_3$).
\end{enumerate}

\end{document}
